%%%%%%%%%%%%%%%%%%%%%%%%%%%%%%%%%%%%%%%%%%%%%%%%%%%%%%%%%%%%
% 6.8 PROBABILISTIC COMBINATORICS
% The Composition of Advanced Mathematics
%%%%%%%%%%%%%%%%%%%%%%%%%%%%%%%%%%%%%%%%%%%%%%%%%%%%%%%%%%%%

\section{Probabilistic Combinatorics}
\label{sec:probabilistic-combinatorics}

\prerequisite{
Basic probability, random variables, expectation, variance, indicator
variables, graph terminology, combinatorial counting, extremal
combinatorics, and familiarity with discrete structures.
}

\learningobjectives{
By the end of this section, the reader should be able to:
\begin{itemize}
    \item explain the probabilistic method;
    \item use positive probability to prove existence;
    \item use indicator random variables in combinatorial arguments;
    \item apply linearity of expectation;
    \item use expectation to prove existence of large or small
          combinatorial structures;
    \item apply the alteration method;
    \item use the first moment method;
    \item use the union bound in probabilistic existence proofs;
    \item apply Markov's inequality to combinatorial random variables;
    \item apply Chebyshev's inequality when variance information is
          available;
    \item understand the second moment method;
    \item recognize the role of independence and limited dependence;
    \item understand the Lov\'asz Local Lemma at an introductory level;
    \item work with the Erd\H{o}s--R\'enyi random graph model \(G(n,p)\);
    \item calculate expected numbers of edges, triangles, and other
          subgraphs;
    \item understand threshold phenomena in random graphs;
    \item distinguish deterministic extremal results from probabilistic
          existence results;
    \item understand random constructions, deletion arguments, and
          randomized selection;
    \item recognize connections among probability, extremal
          combinatorics, graph theory, Ramsey theory, and algorithms.
\end{itemize}
}

%%%%%%%%%%%%%%%%%%%%%%%%%%%%%%%%%%%%%%%%%%%%%%%%%%%%%%%%%%%%
\subsection{What Is Probabilistic Combinatorics?}
\label{subsec:what-is-probabilistic-combinatorics}
%%%%%%%%%%%%%%%%%%%%%%%%%%%%%%%%%%%%%%%%%%%%%%%%%%%%%%%%%%%%

Probabilistic combinatorics studies discrete structures using
probability.

The central idea is often surprising:

\[
\boxed{
\text{randomness can prove deterministic existence}.
}
\]

Suppose a finite probability space consists entirely of combinatorial
objects.

If a randomly chosen object has a desired property with positive
probability, then at least one object with that property must exist.

Symbolically,

\[
\boxed{
\mathbb{P}(A)>0
\Longrightarrow
A\neq\varnothing.
}
\]

This is the foundation of the probabilistic method.

%%%%%%%%%%%%%%%%%%%%%%%%%%%%%%%%%%%%%%%%%%%%%%%%%%%%%%%%%%%%
\subsection{The Probabilistic Method}
\label{subsec:probabilistic-method}
%%%%%%%%%%%%%%%%%%%%%%%%%%%%%%%%%%%%%%%%%%%%%%%%%%%%%%%%%%%%

\begin{definitionbox}{Probabilistic Method}
The probabilistic method proves that a combinatorial object exists by
defining a random object and showing that the desired property occurs
with positive probability.
\end{definitionbox}

The method often proceeds as follows:

\begin{enumerate}
    \item define a random combinatorial object;
    \item define bad events or undesirable statistics;
    \item estimate probabilities or expectations;
    \item show that the bad behavior cannot occur for every object;
    \item conclude that at least one desired object exists.
\end{enumerate}

The proof may not explicitly construct the desired object.

%%%%%%%%%%%%%%%%%%%%%%%%%%%%%%%%%%%%%%%%%%%%%%%%%%%%%%%%%%%%
\subsection{Existence from Expectation}
\label{subsec:existence-from-expectation}
%%%%%%%%%%%%%%%%%%%%%%%%%%%%%%%%%%%%%%%%%%%%%%%%%%%%%%%%%%%%

Suppose \(X\) is a real-valued random variable.

If

\[
\mathbb{E}[X]=\mu,
\]

then at least one outcome satisfies

\[
X\geq\mu
\]

and at least one outcome satisfies

\[
X\leq\mu.
\]

Otherwise every outcome would lie strictly below or strictly above the
average.

Therefore expectation immediately yields deterministic existence
results.

%%%%%%%%%%%%%%%%%%%%%%%%%%%%%%%%%%%%%%%%%%%%%%%%%%%%%%%%%%%%
\subsection{Indicator Random Variables}
\label{subsec:indicator-random-variables}
%%%%%%%%%%%%%%%%%%%%%%%%%%%%%%%%%%%%%%%%%%%%%%%%%%%%%%%%%%%%

For an event \(A\), define the indicator variable

\[
\boxed{
\mathbf{1}_A
=
\begin{cases}
1,&A\text{ occurs},\\
0,&A\text{ does not occur}.
\end{cases}
}
\]

Then

\[
\boxed{
\mathbb{E}[\mathbf{1}_A]
=
\mathbb{P}(A).
}
\]

Indicators convert combinatorial counting questions into sums of random
variables.

If

\[
X
=
\sum_{i=1}^{m}\mathbf{1}_{A_i},
\]

then \(X\) counts how many of the events \(A_i\) occur.

%%%%%%%%%%%%%%%%%%%%%%%%%%%%%%%%%%%%%%%%%%%%%%%%%%%%%%%%%%%%
\subsection{Linearity of Expectation}
\label{subsec:linearity-expectation-combinatorics}
%%%%%%%%%%%%%%%%%%%%%%%%%%%%%%%%%%%%%%%%%%%%%%%%%%%%%%%%%%%%

\begin{theorem}[Linearity of Expectation]
For random variables

\[
X_1,\ldots,X_m,
\]

we have

\[
\boxed{
\mathbb{E}
\left[
\sum_{i=1}^{m}X_i
\right]
=
\sum_{i=1}^{m}\mathbb{E}[X_i].
}
\]
\end{theorem}

Independence is not required.

This makes expectation especially useful in combinatorics because a
complicated random count can often be decomposed into simple indicator
variables.

%%%%%%%%%%%%%%%%%%%%%%%%%%%%%%%%%%%%%%%%%%%%%%%%%%%%%%%%%%%%
\subsection{Worked Example: Expected Fixed Points}
\label{subsec:worked-expected-fixed-points}
%%%%%%%%%%%%%%%%%%%%%%%%%%%%%%%%%%%%%%%%%%%%%%%%%%%%%%%%%%%%

\begin{workedexamplebox}{Fixed Points of a Random Permutation}

Choose a permutation

\[
\pi\in S_n
\]

uniformly at random.

Let

\[
X_i
=
\mathbf{1}_{\{\pi(i)=i\}}.
\]

Then the total number of fixed points is

\[
X
=
\sum_{i=1}^{n}X_i.
\]

For each \(i\),

\[
\mathbb{P}(\pi(i)=i)
=
\frac1n.
\]

Hence,

\[
\mathbb{E}[X_i]
=
\frac1n.
\]

By linearity,

\[
\mathbb{E}[X]
=
\sum_{i=1}^{n}\frac1n
=
1.
\]

\begin{answerbox}
\[
\boxed{
\mathbb{E}[X]=1.
}
\]
\end{answerbox}

Thus a uniformly random permutation has one fixed point on average.

\end{workedexamplebox}

%%%%%%%%%%%%%%%%%%%%%%%%%%%%%%%%%%%%%%%%%%%%%%%%%%%%%%%%%%%%
\subsection{Random Subsets}
\label{subsec:random-subsets}
%%%%%%%%%%%%%%%%%%%%%%%%%%%%%%%%%%%%%%%%%%%%%%%%%%%%%%%%%%%%

A common random construction selects each element independently with
probability \(p\).

Let

\[
S=\{1,\ldots,n\}.
\]

Construct a random subset

\[
R\subseteq S
\]

by including each element independently with probability \(p\).

Then

\[
|R|
=
\sum_{i=1}^{n}X_i,
\]

where

\[
X_i
=
\mathbf{1}_{\{i\in R\}}.
\]

Since

\[
\mathbb{E}[X_i]=p,
\]

we obtain

\[
\boxed{
\mathbb{E}[|R|]=np.
}
\]

%%%%%%%%%%%%%%%%%%%%%%%%%%%%%%%%%%%%%%%%%%%%%%%%%%%%%%%%%%%%
\subsection{The First Moment Method}
\label{subsec:first-moment-method}
%%%%%%%%%%%%%%%%%%%%%%%%%%%%%%%%%%%%%%%%%%%%%%%%%%%%%%%%%%%%

The expectation

\[
\mathbb{E}[X]
\]

is also called the first moment of \(X\).

The first moment method uses expectation to prove existence or
nonexistence.

A particularly useful principle is:

\[
\boxed{
X\geq0
\quad\text{and}\quad
\mathbb{E}[X]<1
\Longrightarrow
\mathbb{P}(X=0)>0
}
\]

when \(X\) is integer-valued.

If every outcome had \(X\geq1\), then

\[
\mathbb{E}[X]\geq1,
\]

a contradiction.

%%%%%%%%%%%%%%%%%%%%%%%%%%%%%%%%%%%%%%%%%%%%%%%%%%%%%%%%%%%%
\subsection{Worked Example: Avoiding Bad Objects}
\label{subsec:worked-first-moment}
%%%%%%%%%%%%%%%%%%%%%%%%%%%%%%%%%%%%%%%%%%%%%%%%%%%%%%%%%%%%

\begin{workedexamplebox}{Existence from Expected Number of Violations}

Suppose a random construction produces a nonnegative integer-valued
random variable \(X\), where \(X\) counts bad configurations.

Assume

\[
\mathbb{E}[X]
=
\frac34.
\]

If every construction contained at least one bad configuration, then

\[
X\geq1
\]

always, so

\[
\mathbb{E}[X]\geq1.
\]

This contradicts

\[
\mathbb{E}[X]=\frac34.
\]

\begin{answerbox}
\[
\boxed{
\text{There exists an outcome with }X=0.
}
\]
\end{answerbox}

\end{workedexamplebox}

%%%%%%%%%%%%%%%%%%%%%%%%%%%%%%%%%%%%%%%%%%%%%%%%%%%%%%%%%%%%
\subsection{Random Graphs}
\label{subsec:random-graphs}
%%%%%%%%%%%%%%%%%%%%%%%%%%%%%%%%%%%%%%%%%%%%%%%%%%%%%%%%%%%%

One of the most important random combinatorial models is the
Erd\H{o}s--R\'enyi random graph.

\begin{definitionbox}{Erd\H{o}s--R\'enyi Random Graph \(G(n,p)\)}
The random graph

\[
G(n,p)
\]

has vertex set

\[
[n]=\{1,\ldots,n\},
\]

and every possible edge is included independently with probability
\(p\).
\end{definitionbox}

There are

\[
\binom{n}{2}
\]

possible edges.

%%%%%%%%%%%%%%%%%%%%%%%%%%%%%%%%%%%%%%%%%%%%%%%%%%%%%%%%%%%%
\subsection{Expected Number of Edges}
\label{subsec:expected-edges-random-graph}
%%%%%%%%%%%%%%%%%%%%%%%%%%%%%%%%%%%%%%%%%%%%%%%%%%%%%%%%%%%%

Let

\[
X
=
|E(G)|.
\]

For every pair

\[
\{i,j\},
\]

define

\[
X_{ij}
=
\mathbf{1}_{\{\{i,j\}\in E(G)\}}.
\]

Then

\[
X
=
\sum_{i<j}X_{ij}.
\]

Since

\[
\mathbb{E}[X_{ij}]
=
p,
\]

linearity gives

\[
\boxed{
\mathbb{E}[|E(G)|]
=
p\binom{n}{2}.
}
\]

%%%%%%%%%%%%%%%%%%%%%%%%%%%%%%%%%%%%%%%%%%%%%%%%%%%%%%%%%%%%
\subsection{Expected Degree}
\label{subsec:expected-degree-random-graph}
%%%%%%%%%%%%%%%%%%%%%%%%%%%%%%%%%%%%%%%%%%%%%%%%%%%%%%%%%%%%

Fix a vertex \(v\).

It has \(n-1\) possible neighbors.

Each edge is present independently with probability \(p\).

Therefore,

\[
\deg(v)
\sim
\operatorname{Bin}(n-1,p).
\]

Hence,

\[
\boxed{
\mathbb{E}[\deg(v)]
=
(n-1)p.
}
\]

The notation

\[
\operatorname{Bin}(n-1,p)
\]

denotes a binomial distribution.

%%%%%%%%%%%%%%%%%%%%%%%%%%%%%%%%%%%%%%%%%%%%%%%%%%%%%%%%%%%%
\subsection{Expected Number of Triangles}
\label{subsec:expected-triangles}
%%%%%%%%%%%%%%%%%%%%%%%%%%%%%%%%%%%%%%%%%%%%%%%%%%%%%%%%%%%%

There are

\[
\binom{n}{3}
\]

sets of three vertices.

For each triple \(S\), let

\[
X_S
\]

be the indicator that those vertices form a triangle.

Three edges must be present.

Therefore,

\[
\mathbb{P}(X_S=1)
=
p^3.
\]

If \(X\) is the total number of triangles,

\[
X
=
\sum_{S\in\binom{[n]}3}X_S.
\]

Thus,

\[
\boxed{
\mathbb{E}[X]
=
\binom{n}{3}p^3.
}
\]

%%%%%%%%%%%%%%%%%%%%%%%%%%%%%%%%%%%%%%%%%%%%%%%%%%%%%%%%%%%%
\subsection{Worked Example: Random Graph Triangles}
\label{subsec:worked-random-triangles}
%%%%%%%%%%%%%%%%%%%%%%%%%%%%%%%%%%%%%%%%%%%%%%%%%%%%%%%%%%%%

\begin{workedexamplebox}{Expected Triangles in \(G(10,\tfrac12)\)}

For

\[
n=10,
\qquad
p=\frac12,
\]

the expected number of triangles is

\[
\binom{10}{3}
\left(\frac12\right)^3.
\]

Since

\[
\binom{10}{3}=120,
\]

we obtain

\[
120\cdot\frac18=15.
\]

\begin{answerbox}
\[
\boxed{
\mathbb{E}[X]=15.
}
\]
\end{answerbox}

\end{workedexamplebox}

%%%%%%%%%%%%%%%%%%%%%%%%%%%%%%%%%%%%%%%%%%%%%%%%%%%%%%%%%%%%
\subsection{Expected Number of Copies of a Fixed Graph}
\label{subsec:expected-fixed-subgraph}
%%%%%%%%%%%%%%%%%%%%%%%%%%%%%%%%%%%%%%%%%%%%%%%%%%%%%%%%%%%%

Let \(H\) be a fixed graph with

\[
v(H)
\]

vertices and

\[
e(H)
\]

edges.

A specified labeled copy of \(H\) requires all of its \(e(H)\) edges to
appear.

Therefore its probability is

\[
p^{e(H)}.
\]

Thus the expected number of copies of \(H\) typically has the form

\[
\boxed{
\text{number of candidate copies}
\times
p^{e(H)}.
}
\]

This simple principle is extremely useful in random graph theory.

%%%%%%%%%%%%%%%%%%%%%%%%%%%%%%%%%%%%%%%%%%%%%%%%%%%%%%%%%%%%
\subsection{Random Colorings}
\label{subsec:random-colorings}
%%%%%%%%%%%%%%%%%%%%%%%%%%%%%%%%%%%%%%%%%%%%%%%%%%%%%%%%%%%%

Another common model assigns colors randomly.

Suppose each object receives one of \(q\) colors independently and
uniformly.

Then each color has probability

\[
\frac1q.
\]

A specified set of \(r\) objects is monochromatic with probability

\[
\boxed{
q\left(\frac1q\right)^r
=
q^{1-r}.
}
\]

Random coloring arguments are fundamental in Ramsey theory.

%%%%%%%%%%%%%%%%%%%%%%%%%%%%%%%%%%%%%%%%%%%%%%%%%%%%%%%%%%%%
\subsection{Worked Example: Monochromatic Triples}
\label{subsec:worked-random-coloring}
%%%%%%%%%%%%%%%%%%%%%%%%%%%%%%%%%%%%%%%%%%%%%%%%%%%%%%%%%%%%

\begin{workedexamplebox}{A Random Two-Coloring}

Suppose three specified objects are independently colored red or blue
with equal probability.

The probability that all three receive the same color is

\[
2\left(\frac12\right)^3
=
\frac14.
\]

\begin{answerbox}
\[
\boxed{
\mathbb{P}(\text{monochromatic})=\frac14.
}
\]
\end{answerbox}

\end{workedexamplebox}

%%%%%%%%%%%%%%%%%%%%%%%%%%%%%%%%%%%%%%%%%%%%%%%%%%%%%%%%%%%%
\subsection{The Union Bound}
\label{subsec:union-bound}
%%%%%%%%%%%%%%%%%%%%%%%%%%%%%%%%%%%%%%%%%%%%%%%%%%%%%%%%%%%%

\begin{theorem}[Union Bound]
For events

\[
A_1,\ldots,A_m,
\]

we have

\[
\boxed{
\mathbb{P}
\left(
\bigcup_{i=1}^{m}A_i
\right)
\leq
\sum_{i=1}^{m}\mathbb{P}(A_i).
}
\]
\end{theorem}

Independence is not required.

The union bound is one of the most frequently used tools in
probabilistic combinatorics.

%%%%%%%%%%%%%%%%%%%%%%%%%%%%%%%%%%%%%%%%%%%%%%%%%%%%%%%%%%%%
\subsection{Existence from the Union Bound}
\label{subsec:existence-union-bound}
%%%%%%%%%%%%%%%%%%%%%%%%%%%%%%%%%%%%%%%%%%%%%%%%%%%%%%%%%%%%

Suppose \(A_1,\ldots,A_m\) are all undesirable events.

If

\[
\boxed{
\sum_{i=1}^{m}\mathbb{P}(A_i)<1,
}
\]

then

\[
\mathbb{P}
\left(
\bigcup_iA_i
\right)
<1.
\]

Therefore,

\[
\mathbb{P}
\left(
\bigcap_iA_i^c
\right)
>0.
\]

Hence there exists an outcome in which none of the bad events occur.

%%%%%%%%%%%%%%%%%%%%%%%%%%%%%%%%%%%%%%%%%%%%%%%%%%%%%%%%%%%%
\subsection{Union Bound Versus Inclusion--Exclusion}
\label{subsec:union-bound-inclusion-exclusion}
%%%%%%%%%%%%%%%%%%%%%%%%%%%%%%%%%%%%%%%%%%%%%%%%%%%%%%%%%%%%

Inclusion--exclusion gives an exact formula for

\[
\mathbb{P}\left(\bigcup_iA_i\right),
\]

but may require many intersection terms.

The union bound gives only an upper bound:

\[
\mathbb{P}\left(\bigcup_iA_i\right)
\leq
\sum_i\mathbb{P}(A_i).
\]

Its advantage is simplicity.

For existence proofs, exact probability is often unnecessary.

It is enough to show the total bad probability is less than \(1\).

%%%%%%%%%%%%%%%%%%%%%%%%%%%%%%%%%%%%%%%%%%%%%%%%%%%%%%%%%%%%
\subsection{Ramsey Lower Bounds by Random Coloring}
\label{subsec:ramsey-random-coloring-preview}
%%%%%%%%%%%%%%%%%%%%%%%%%%%%%%%%%%%%%%%%%%%%%%%%%%%%%%%%%%%%

Consider coloring the edges of

\[
K_n
\]

red or blue independently with equal probability.

For a fixed set of \(k\) vertices, the probability that all

\[
\binom{k}{2}
\]

edges receive the same color is

\[
\boxed{
2^{1-\binom{k}{2}}.
}
\]

There are

\[
\binom{n}{k}
\]

possible \(k\)-vertex subsets.

Therefore, by the union bound, the probability of having at least one
monochromatic \(K_k\) is at most

\[
\boxed{
\binom{n}{k}
2^{1-\binom{k}{2}}.
}
\]

If this quantity is less than \(1\), then a coloring with no
monochromatic \(K_k\) exists.

This yields lower bounds for Ramsey numbers.

%%%%%%%%%%%%%%%%%%%%%%%%%%%%%%%%%%%%%%%%%%%%%%%%%%%%%%%%%%%%
\subsection{Worked Example: Ramsey-Style Existence}
\label{subsec:worked-ramsey-probabilistic}
%%%%%%%%%%%%%%%%%%%%%%%%%%%%%%%%%%%%%%%%%%%%%%%%%%%%%%%%%%%%

\begin{workedexamplebox}{Avoiding a Monochromatic \(K_4\)}

Color the edges of \(K_n\) independently red or blue.

For any fixed \(4\)-vertex set, there are

\[
\binom42=6
\]

edges.

The probability all six are the same color is

\[
2\left(\frac12\right)^6
=
\frac1{32}.
\]

There are

\[
\binom{n}{4}
\]

candidate \(K_4\)'s.

Therefore,

\[
\mathbb{P}(\text{some monochromatic }K_4)
\leq
\frac{\binom{n}{4}}{32}.
\]

Whenever

\[
\frac{\binom{n}{4}}{32}<1,
\]

a coloring avoiding monochromatic \(K_4\)'s must exist.

\begin{answerbox}
\[
\boxed{
\binom{n}{4}<32
\Longrightarrow
\text{such a coloring exists}.
}
\]
\end{answerbox}

\end{workedexamplebox}

%%%%%%%%%%%%%%%%%%%%%%%%%%%%%%%%%%%%%%%%%%%%%%%%%%%%%%%%%%%%
\subsection{Markov's Inequality}
\label{subsec:markovs-inequality-combinatorics}
%%%%%%%%%%%%%%%%%%%%%%%%%%%%%%%%%%%%%%%%%%%%%%%%%%%%%%%%%%%%

\begin{theorem}[Markov's Inequality]
If \(X\) is a nonnegative random variable and \(a>0\), then

\[
\boxed{
\mathbb{P}(X\geq a)
\leq
\frac{\mathbb{E}[X]}{a}.
}
\]
\end{theorem}

Markov's inequality converts an expectation bound into a tail
probability bound.

It requires no independence assumption.

%%%%%%%%%%%%%%%%%%%%%%%%%%%%%%%%%%%%%%%%%%%%%%%%%%%%%%%%%%%%
\subsection{Worked Example: Markov Bound}
\label{subsec:worked-markov}
%%%%%%%%%%%%%%%%%%%%%%%%%%%%%%%%%%%%%%%%%%%%%%%%%%%%%%%%%%%%

\begin{workedexamplebox}{Bounding a Large Count}

Suppose \(X\geq0\) and

\[
\mathbb{E}[X]=10.
\]

Then

\[
\mathbb{P}(X\geq50)
\leq
\frac{10}{50}
=
\frac15.
\]

\begin{answerbox}
\[
\boxed{
\mathbb{P}(X\geq50)\leq\frac15.
}
\]
\end{answerbox}

\end{workedexamplebox}

%%%%%%%%%%%%%%%%%%%%%%%%%%%%%%%%%%%%%%%%%%%%%%%%%%%%%%%%%%%%
\subsection{Variance}
\label{subsec:variance-probabilistic-combinatorics}
%%%%%%%%%%%%%%%%%%%%%%%%%%%%%%%%%%%%%%%%%%%%%%%%%%%%%%%%%%%%

Expectation measures average value.

Variance measures spread around the mean.

For a random variable \(X\),

\[
\boxed{
\operatorname{Var}(X)
=
\mathbb{E}
\left[
(X-\mathbb{E}[X])^2
\right].
}
\]

Equivalently,

\[
\boxed{
\operatorname{Var}(X)
=
\mathbb{E}[X^2]
-
\mathbb{E}[X]^2.
}
\]

Small variance indicates that \(X\) is concentrated near its mean.

%%%%%%%%%%%%%%%%%%%%%%%%%%%%%%%%%%%%%%%%%%%%%%%%%%%%%%%%%%%%
\subsection{Variance of a Sum}
\label{subsec:variance-sum}
%%%%%%%%%%%%%%%%%%%%%%%%%%%%%%%%%%%%%%%%%%%%%%%%%%%%%%%%%%%%

For random variables \(X_1,\ldots,X_n\),

\[
\operatorname{Var}
\left(
\sum_iX_i
\right)
=
\sum_i\operatorname{Var}(X_i)
+
2\sum_{i<j}
\operatorname{Cov}(X_i,X_j).
\]

If the variables are independent, all covariance terms vanish.

Then

\[
\boxed{
\operatorname{Var}
\left(
\sum_iX_i
\right)
=
\sum_i\operatorname{Var}(X_i).
}
\]

%%%%%%%%%%%%%%%%%%%%%%%%%%%%%%%%%%%%%%%%%%%%%%%%%%%%%%%%%%%%
\subsection{Chebyshev's Inequality}
\label{subsec:chebyshev-combinatorics}
%%%%%%%%%%%%%%%%%%%%%%%%%%%%%%%%%%%%%%%%%%%%%%%%%%%%%%%%%%%%

\begin{theorem}[Chebyshev's Inequality]
For a random variable \(X\) with finite variance and every \(t>0\),

\[
\boxed{
\mathbb{P}
\left(
|X-\mathbb{E}[X]|\geq t
\right)
\leq
\frac{\operatorname{Var}(X)}{t^2}.
}
\]
\end{theorem}

This provides concentration around the expectation.

%%%%%%%%%%%%%%%%%%%%%%%%%%%%%%%%%%%%%%%%%%%%%%%%%%%%%%%%%%%%
\subsection{Concentration of Random Graph Edges}
\label{subsec:random-graph-edge-concentration}
%%%%%%%%%%%%%%%%%%%%%%%%%%%%%%%%%%%%%%%%%%%%%%%%%%%%%%%%%%%%

For

\[
G\sim G(n,p),
\]

the number of edges satisfies

\[
X
\sim
\operatorname{Bin}
\left(
\binom n2,p
\right).
\]

Thus,

\[
\mathbb{E}[X]
=
p\binom n2
\]

and

\[
\operatorname{Var}(X)
=
p(1-p)\binom n2.
\]

Chebyshev's inequality shows that \(X\) becomes increasingly
concentrated around its mean as the graph grows.

%%%%%%%%%%%%%%%%%%%%%%%%%%%%%%%%%%%%%%%%%%%%%%%%%%%%%%%%%%%%
\subsection{Second Moment Method}
\label{subsec:second-moment-method}
%%%%%%%%%%%%%%%%%%%%%%%%%%%%%%%%%%%%%%%%%%%%%%%%%%%%%%%%%%%%

The first moment method is often used to show that an undesirable count
is zero.

The second moment method is frequently used to show that a desired
count is positive.

A useful inequality is

\[
\boxed{
\mathbb{P}(X>0)
\geq
\frac{\mathbb{E}[X]^2}{\mathbb{E}[X^2]}
}
\]

for nonnegative random variables \(X\).

This is a form of the second moment method.

%%%%%%%%%%%%%%%%%%%%%%%%%%%%%%%%%%%%%%%%%%%%%%%%%%%%%%%%%%%%
\subsection{Why the Second Moment Helps}
\label{subsec:why-second-moment}
%%%%%%%%%%%%%%%%%%%%%%%%%%%%%%%%%%%%%%%%%%%%%%%%%%%%%%%%%%%%

A large expectation alone does not guarantee that

\[
X>0
\]

with high probability.

A random variable may have a large average caused by rare, enormous
values.

The second moment measures the degree of fluctuation.

If

\[
\mathbb{E}[X^2]
\]

is not much larger than

\[
\mathbb{E}[X]^2,
\]

then \(X\) cannot usually be zero.

%%%%%%%%%%%%%%%%%%%%%%%%%%%%%%%%%%%%%%%%%%%%%%%%%%%%%%%%%%%%
\subsection{Independent Events}
\label{subsec:independent-events-combinatorics}
%%%%%%%%%%%%%%%%%%%%%%%%%%%%%%%%%%%%%%%%%%%%%%%%%%%%%%%%%%%%

Events \(A\) and \(B\) are independent when

\[
\boxed{
\mathbb{P}(A\cap B)
=
\mathbb{P}(A)\mathbb{P}(B).
}
\]

Mutual independence of several events requires the corresponding
identity for every subcollection.

Many random combinatorial models are built from independent elementary
choices.

For example, the edges in \(G(n,p)\) are independent.

However, larger graph events are usually not independent.

%%%%%%%%%%%%%%%%%%%%%%%%%%%%%%%%%%%%%%%%%%%%%%%%%%%%%%%%%%%%
\subsection{Dependence Between Subgraph Events}
\label{subsec:subgraph-event-dependence}
%%%%%%%%%%%%%%%%%%%%%%%%%%%%%%%%%%%%%%%%%%%%%%%%%%%%%%%%%%%%

Consider two potential triangles in \(G(n,p)\).

If they have no edges in common, their edge requirements are
independent.

If they share an edge, the events are dependent.

Thus local overlap creates dependence.

Understanding these dependency patterns is important for:

\begin{itemize}
    \item variance calculations;
    \item second moment arguments;
    \item concentration inequalities;
    \item the Lov\'asz Local Lemma.
\end{itemize}

%%%%%%%%%%%%%%%%%%%%%%%%%%%%%%%%%%%%%%%%%%%%%%%%%%%%%%%%%%%%
\subsection{The Lov\'asz Local Lemma}
\label{subsec:lovasz-local-lemma}
%%%%%%%%%%%%%%%%%%%%%%%%%%%%%%%%%%%%%%%%%%%%%%%%%%%%%%%%%%%%

The union bound may fail when there are many bad events, even when each
event is individually unlikely.

The Lov\'asz Local Lemma shows that if each bad event is rare and
depends on only a limited number of other bad events, then all bad
events may simultaneously be avoided.

%%%%%%%%%%%%%%%%%%%%%%%%%%%%%%%%%%%%%%%%%%%%%%%%%%%%%%%%%%%%
\subsection{Symmetric Lov\'asz Local Lemma}
\label{subsec:symmetric-lovasz-local-lemma}
%%%%%%%%%%%%%%%%%%%%%%%%%%%%%%%%%%%%%%%%%%%%%%%%%%%%%%%%%%%%

\begin{theorem}[Symmetric Lov\'asz Local Lemma]
Suppose events

\[
A_1,\ldots,A_m
\]

satisfy

\[
\mathbb{P}(A_i)\leq p
\]

for every \(i\), and each event is mutually independent of all but at
most \(d\) other events.

If

\[
\boxed{
e p(d+1)\leq1,
}
\]

then

\[
\boxed{
\mathbb{P}
\left(
\bigcap_{i=1}^{m}A_i^c
\right)
>0.
}
\]
\end{theorem}

Here,

\[
e
\]

is Euler's number.

%%%%%%%%%%%%%%%%%%%%%%%%%%%%%%%%%%%%%%%%%%%%%%%%%%%%%%%%%%%%
\subsection{Why the Local Lemma Is Stronger Than the Union Bound}
\label{subsec:local-lemma-versus-union-bound}
%%%%%%%%%%%%%%%%%%%%%%%%%%%%%%%%%%%%%%%%%%%%%%%%%%%%%%%%%%%%

The union bound depends on the total number of bad events:

\[
\sum_i\mathbb{P}(A_i).
\]

The Local Lemma instead depends primarily on:

\[
\boxed{
\text{event probability}
+
\text{local dependency}.
}
\]

Therefore, even if there are enormously many bad events, all may still
be avoidable if each event interacts with only a small neighborhood of
others.

%%%%%%%%%%%%%%%%%%%%%%%%%%%%%%%%%%%%%%%%%%%%%%%%%%%%%%%%%%%%
\subsection{Randomized Coloring}
\label{subsec:randomized-coloring}
%%%%%%%%%%%%%%%%%%%%%%%%%%%%%%%%%%%%%%%%%%%%%%%%%%%%%%%%%%%%

Suppose vertices of a graph are assigned colors independently.

A bad event might be

\[
A_e
=
\{\text{edge }e\text{ has equal colors at both endpoints}\}.
\]

Each \(A_e\) depends only on events involving neighboring edges.

This local dependence structure makes coloring problems natural
applications of the Local Lemma.

%%%%%%%%%%%%%%%%%%%%%%%%%%%%%%%%%%%%%%%%%%%%%%%%%%%%%%%%%%%%
\subsection{Alteration Method}
\label{subsec:alteration-method-probabilistic}
%%%%%%%%%%%%%%%%%%%%%%%%%%%%%%%%%%%%%%%%%%%%%%%%%%%%%%%%%%%%

The alteration method first constructs a random object that is nearly
good.

Then a small number of modifications remove the remaining bad
configurations.

The structure is

\[
\boxed{
\text{random construction}
\longrightarrow
\text{identify violations}
\longrightarrow
\text{repair violations}.
}
\]

This often produces better bounds than requiring a random object to be
perfect immediately.

%%%%%%%%%%%%%%%%%%%%%%%%%%%%%%%%%%%%%%%%%%%%%%%%%%%%%%%%%%%%
\subsection{Independent Sets by Alteration}
\label{subsec:independent-set-alteration}
%%%%%%%%%%%%%%%%%%%%%%%%%%%%%%%%%%%%%%%%%%%%%%%%%%%%%%%%%%%%

Let

\[
G=(V,E)
\]

have

\[
|V|=n,
\qquad
|E|=m.
\]

Choose every vertex independently with probability \(p\).

Let \(X\) be the number of selected vertices and \(Y\) the number of
edges whose endpoints are both selected.

Then

\[
\mathbb{E}[X]=pn
\]

and

\[
\mathbb{E}[Y]=p^2m.
\]

Delete one endpoint from every selected edge.

The remaining set is independent and has size at least

\[
X-Y.
\]

Therefore,

\[
\boxed{
\alpha(G)
\geq
pn-p^2m
}
\]

for some choice of the random subset.

%%%%%%%%%%%%%%%%%%%%%%%%%%%%%%%%%%%%%%%%%%%%%%%%%%%%%%%%%%%%
\subsection{Optimizing an Alteration Bound}
\label{subsec:optimizing-alteration}
%%%%%%%%%%%%%%%%%%%%%%%%%%%%%%%%%%%%%%%%%%%%%%%%%%%%%%%%%%%%

Consider

\[
f(p)
=
pn-p^2m.
\]

Differentiate:

\[
f'(p)
=
n-2mp.
\]

The unconstrained maximum occurs at

\[
\boxed{
p=\frac{n}{2m}
}
\]

when this lies in \([0,1]\).

Substituting yields

\[
f(p)
=
\frac{n^2}{4m}.
\]

Thus, in the appropriate range,

\[
\boxed{
\alpha(G)
\geq
\frac{n^2}{4m}.
}
\]

%%%%%%%%%%%%%%%%%%%%%%%%%%%%%%%%%%%%%%%%%%%%%%%%%%%%%%%%%%%%
\subsection{Random Sampling and Deletion}
\label{subsec:random-sampling-deletion}
%%%%%%%%%%%%%%%%%%%%%%%%%%%%%%%%%%%%%%%%%%%%%%%%%%%%%%%%%%%%

More generally, suppose a random sample contains:

\[
X
=
\text{number of useful objects}
\]

and

\[
Y
=
\text{number of bad configurations}.
\]

If one deletion can destroy each bad configuration, then the repaired
object has size at least

\[
X-Y.
\]

Hence,

\[
\boxed{
\mathbb{E}[X-Y]
=
\mathbb{E}[X]-\mathbb{E}[Y].
}
\]

A positive lower bound on this expectation guarantees a deterministic
object at least that large.

%%%%%%%%%%%%%%%%%%%%%%%%%%%%%%%%%%%%%%%%%%%%%%%%%%%%%%%%%%%%
\subsection{Random Graph Independence Number}
\label{subsec:random-graph-independence-preview}
%%%%%%%%%%%%%%%%%%%%%%%%%%%%%%%%%%%%%%%%%%%%%%%%%%%%%%%%%%%%

In a random graph \(G(n,p)\), an independent set of size \(k\) is a
\(k\)-vertex set containing no edges.

For a fixed \(k\)-set, the probability of independence is

\[
\boxed{
(1-p)^{\binom{k}{2}}.
}
\]

There are

\[
\binom nk
\]

candidate sets.

Thus the expected number of independent \(k\)-sets is

\[
\boxed{
\binom nk
(1-p)^{\binom{k}{2}}.
}
\]

This first-moment expression can be used to estimate the independence
number of random graphs.

%%%%%%%%%%%%%%%%%%%%%%%%%%%%%%%%%%%%%%%%%%%%%%%%%%%%%%%%%%%%
\subsection{Random Graph Cliques}
\label{subsec:random-graph-cliques}
%%%%%%%%%%%%%%%%%%%%%%%%%%%%%%%%%%%%%%%%%%%%%%%%%%%%%%%%%%%%

A fixed \(k\)-vertex set forms a clique when all

\[
\binom{k}{2}
\]

possible edges are present.

Thus,

\[
\boxed{
\mathbb{P}
(\text{fixed }k\text{-set is a clique})
=
p^{\binom{k}{2}}.
}
\]

The expected number of \(K_k\)'s is

\[
\boxed{
\binom nk
p^{\binom{k}{2}}.
}
\]

This expectation helps identify the scale at which cliques begin to
appear.

%%%%%%%%%%%%%%%%%%%%%%%%%%%%%%%%%%%%%%%%%%%%%%%%%%%%%%%%%%%%
\subsection{Threshold Phenomena}
\label{subsec:threshold-phenomena}
%%%%%%%%%%%%%%%%%%%%%%%%%%%%%%%%%%%%%%%%%%%%%%%%%%%%%%%%%%%%

Many random graph properties change rapidly as \(p\) passes a critical
scale.

For a property \(\mathcal{P}\), a function

\[
p_c(n)
\]

may act as a threshold if:

\[
p\ll p_c(n)
\]

typically implies the property is absent, while

\[
p\gg p_c(n)
\]

typically implies the property is present.

The notation

\[
f(n)\ll g(n)
\]

means

\[
\frac{f(n)}{g(n)}\to0.
\]

%%%%%%%%%%%%%%%%%%%%%%%%%%%%%%%%%%%%%%%%%%%%%%%%%%%%%%%%%%%%
\subsection{Threshold for a Single Edge}
\label{subsec:threshold-single-edge}
%%%%%%%%%%%%%%%%%%%%%%%%%%%%%%%%%%%%%%%%%%%%%%%%%%%%%%%%%%%%

For \(G(n,p)\), the expected number of edges is

\[
p\binom n2.
\]

If

\[
p\ll n^{-2},
\]

then this expectation tends to \(0\).

Thus the graph typically has no edges.

If

\[
p\gg n^{-2},
\]

edges begin to appear.

This illustrates the general idea that first moments often suggest
threshold scales.

%%%%%%%%%%%%%%%%%%%%%%%%%%%%%%%%%%%%%%%%%%%%%%%%%%%%%%%%%%%%
\subsection{Threshold for Triangles}
\label{subsec:threshold-triangles-preview}
%%%%%%%%%%%%%%%%%%%%%%%%%%%%%%%%%%%%%%%%%%%%%%%%%%%%%%%%%%%%

The expected number of triangles is

\[
\binom n3p^3.
\]

This becomes order \(1\) when approximately

\[
n^3p^3
\asymp
1.
\]

Hence,

\[
\boxed{
p\asymp n^{-1}
}
\]

is the natural scale for triangle appearance.

A rigorous threshold theorem requires more than expectation alone, but
the first moment identifies the correct scale.

%%%%%%%%%%%%%%%%%%%%%%%%%%%%%%%%%%%%%%%%%%%%%%%%%%%%%%%%%%%%
\subsection{With High Probability}
\label{subsec:with-high-probability}
%%%%%%%%%%%%%%%%%%%%%%%%%%%%%%%%%%%%%%%%%%%%%%%%%%%%%%%%%%%%

An event \(A_n\) occurs \emph{with high probability}, abbreviated
w.h.p., if

\[
\boxed{
\mathbb{P}(A_n)\to1
}
\]

as

\[
n\to\infty.
\]

Random graph theory frequently studies properties that occur with high
probability.

Examples include:

\begin{itemize}
    \item connectivity;
    \item existence of cycles;
    \item clique size;
    \item chromatic number;
    \item degree concentration.
\end{itemize}

%%%%%%%%%%%%%%%%%%%%%%%%%%%%%%%%%%%%%%%%%%%%%%%%%%%%%%%%%%%%
\subsection{Random Variables Counting Structures}
\label{subsec:random-counting-variables}
%%%%%%%%%%%%%%%%%%%%%%%%%%%%%%%%%%%%%%%%%%%%%%%%%%%%%%%%%%%%

A standard probabilistic-combinatorial technique is:

\begin{enumerate}
    \item identify all candidate structures;
    \item assign one indicator variable to each candidate;
    \item sum the indicators;
    \item apply linearity of expectation.
\end{enumerate}

For example,

\[
X
=
\sum_{S}
\mathbf{1}_{\{S\text{ forms the desired structure}\}}.
\]

Then

\[
\mathbb{E}[X]
=
\sum_S
\mathbb{P}(S\text{ forms the desired structure}).
\]

%%%%%%%%%%%%%%%%%%%%%%%%%%%%%%%%%%%%%%%%%%%%%%%%%%%%%%%%%%%%
\subsection{Worked Example: Expected \(k\)-Cliques}
\label{subsec:worked-expected-k-cliques}
%%%%%%%%%%%%%%%%%%%%%%%%%%%%%%%%%%%%%%%%%%%%%%%%%%%%%%%%%%%%

\begin{workedexamplebox}{Counting Expected Cliques}

Let

\[
G\sim G(n,p).
\]

For every

\[
S\in\binom{[n]}k,
\]

define

\[
X_S
=
\mathbf{1}_{\{S\text{ induces a }K_k\}}.
\]

Then

\[
X
=
\sum_{S\in\binom{[n]}k}X_S.
\]

A \(K_k\) requires

\[
\binom{k}{2}
\]

edges.

Therefore,

\[
\mathbb{E}[X_S]
=
p^{\binom{k}{2}}.
\]

By linearity,

\begin{answerbox}
\[
\boxed{
\mathbb{E}[X]
=
\binom nk
p^{\binom{k}{2}}.
}
\]
\end{answerbox}

\end{workedexamplebox}

%%%%%%%%%%%%%%%%%%%%%%%%%%%%%%%%%%%%%%%%%%%%%%%%%%%%%%%%%%%%
\subsection{Random Hypergraphs}
\label{subsec:random-hypergraphs}
%%%%%%%%%%%%%%%%%%%%%%%%%%%%%%%%%%%%%%%%%%%%%%%%%%%%%%%%%%%%

The random graph model generalizes to hypergraphs.

For a \(k\)-uniform random hypergraph, every \(k\)-element subset of the
vertex set is included independently with probability \(p\).

This model is often denoted

\[
\boxed{
H^{(k)}(n,p).
}
\]

There are

\[
\binom nk
\]

possible hyperedges.

Hence the expected number of hyperedges is

\[
\boxed{
p\binom nk.
}
\]

%%%%%%%%%%%%%%%%%%%%%%%%%%%%%%%%%%%%%%%%%%%%%%%%%%%%%%%%%%%%
\subsection{Random Set Systems}
\label{subsec:random-set-systems}
%%%%%%%%%%%%%%%%%%%%%%%%%%%%%%%%%%%%%%%%%%%%%%%%%%%%%%%%%%%%

One may also select members of

\[
\mathcal{P}([n])
\]

randomly.

For example, include every subset independently with probability \(p\).

Such models allow probabilistic study of:

\begin{itemize}
    \item intersecting families;
    \item antichains;
    \item covering families;
    \item random simplicial complexes;
    \item random Boolean functions.
\end{itemize}

%%%%%%%%%%%%%%%%%%%%%%%%%%%%%%%%%%%%%%%%%%%%%%%%%%%%%%%%%%%%
\subsection{Random Permutations}
\label{subsec:random-permutations}
%%%%%%%%%%%%%%%%%%%%%%%%%%%%%%%%%%%%%%%%%%%%%%%%%%%%%%%%%%%%

A uniformly random permutation selects each element of

\[
S_n
\]

with probability

\[
\frac1{n!}.
\]

Important statistics include:

\[
\operatorname{fix}(\pi),
\qquad
\operatorname{cyc}(\pi),
\qquad
\operatorname{inv}(\pi),
\qquad
\operatorname{des}(\pi).
\]

Probabilistic methods ask about their expectations, variances, and
limiting distributions.

%%%%%%%%%%%%%%%%%%%%%%%%%%%%%%%%%%%%%%%%%%%%%%%%%%%%%%%%%%%%
\subsection{Expected Number of Inversions}
\label{subsec:expected-inversions}
%%%%%%%%%%%%%%%%%%%%%%%%%%%%%%%%%%%%%%%%%%%%%%%%%%%%%%%%%%%%

Let \(\pi\) be a uniformly random permutation of \(n\) elements.

For each pair

\[
i<j,
\]

define

\[
X_{ij}
=
\mathbf{1}_{\{\pi(i)>\pi(j)\}}.
\]

Exactly one of

\[
\pi(i)>\pi(j)
\]

and

\[
\pi(i)<\pi(j)
\]

occurs, symmetrically.

Therefore,

\[
\mathbb{P}(X_{ij}=1)
=
\frac12.
\]

Since there are

\[
\binom n2
\]

pairs,

\[
\boxed{
\mathbb{E}[\operatorname{inv}(\pi)]
=
\frac12\binom n2
=
\frac{n(n-1)}4.
}
\]

%%%%%%%%%%%%%%%%%%%%%%%%%%%%%%%%%%%%%%%%%%%%%%%%%%%%%%%%%%%%
\subsection{Worked Example: Expected Inversions}
\label{subsec:worked-expected-inversions}
%%%%%%%%%%%%%%%%%%%%%%%%%%%%%%%%%%%%%%%%%%%%%%%%%%%%%%%%%%%%

\begin{workedexamplebox}{Random Permutation of Six Elements}

For

\[
n=6,
\]

the expected number of inversions is

\[
\frac12\binom62.
\]

Since

\[
\binom62=15,
\]

we obtain

\begin{answerbox}
\[
\boxed{
\mathbb{E}[\operatorname{inv}(\pi)]
=
\frac{15}{2}.
}
\]
\end{answerbox}

\end{workedexamplebox}

%%%%%%%%%%%%%%%%%%%%%%%%%%%%%%%%%%%%%%%%%%%%%%%%%%%%%%%%%%%%
\subsection{Expected Number of Cycles}
\label{subsec:expected-cycles-random-permutation}
%%%%%%%%%%%%%%%%%%%%%%%%%%%%%%%%%%%%%%%%%%%%%%%%%%%%%%%%%%%%

For a uniformly random permutation of \(n\) elements, the expected
number of cycles is

\[
\boxed{
H_n
=
1+\frac12+\frac13+\cdots+\frac1n.
}
\]

The quantity \(H_n\) is the \(n\)-th harmonic number.

Asymptotically,

\[
\boxed{
H_n
=
\log n+\gamma+o(1),
}
\]

where

\[
\gamma
\]

is the Euler--Mascheroni constant.

%%%%%%%%%%%%%%%%%%%%%%%%%%%%%%%%%%%%%%%%%%%%%%%%%%%%%%%%%%%%
\subsection{Randomized Algorithms}
\label{subsec:randomized-algorithms-preview}
%%%%%%%%%%%%%%%%%%%%%%%%%%%%%%%%%%%%%%%%%%%%%%%%%%%%%%%%%%%%

Probability may be used not only in proofs but also in algorithms.

A randomized algorithm makes random choices during execution.

Randomization can help:

\begin{itemize}
    \item simplify algorithms;
    \item avoid adversarial inputs;
    \item estimate large quantities;
    \item find structures efficiently;
    \item approximate difficult optimization problems.
\end{itemize}

Probabilistic combinatorics often supplies the mathematical analysis
behind such algorithms.

%%%%%%%%%%%%%%%%%%%%%%%%%%%%%%%%%%%%%%%%%%%%%%%%%%%%%%%%%%%%
\subsection{Monte Carlo and Las Vegas Algorithms}
\label{subsec:monte-carlo-las-vegas}
%%%%%%%%%%%%%%%%%%%%%%%%%%%%%%%%%%%%%%%%%%%%%%%%%%%%%%%%%%%%

A Monte Carlo algorithm may return an incorrect or approximate answer
with controlled probability.

A Las Vegas algorithm always returns a correct answer, but its running
time is random.

These two paradigms illustrate different ways randomness can enter
computation.

%%%%%%%%%%%%%%%%%%%%%%%%%%%%%%%%%%%%%%%%%%%%%%%%%%%%%%%%%%%%
\subsection{Conditional Expectation}
\label{subsec:conditional-expectation-method}
%%%%%%%%%%%%%%%%%%%%%%%%%%%%%%%%%%%%%%%%%%%%%%%%%%%%%%%%%%%%

A probabilistic existence proof can sometimes be converted into a
deterministic algorithm.

Suppose a random construction has objective \(X\) with

\[
\mathbb{E}[X]\geq M.
\]

Make random choices one at a time.

At each step, choose an option that does not decrease the conditional
expectation below its current value.

Eventually all choices become deterministic.

The final object satisfies

\[
\boxed{
X\geq M.
}
\]

This is the method of conditional expectation.

%%%%%%%%%%%%%%%%%%%%%%%%%%%%%%%%%%%%%%%%%%%%%%%%%%%%%%%%%%%%
\subsection{Derandomization}
\label{subsec:derandomization}
%%%%%%%%%%%%%%%%%%%%%%%%%%%%%%%%%%%%%%%%%%%%%%%%%%%%%%%%%%%%

Derandomization replaces a probabilistic construction or randomized
algorithm with a deterministic procedure having comparable guarantees.

Methods include:

\begin{itemize}
    \item conditional expectation;
    \item limited independence;
    \item pseudorandom generators;
    \item discrepancy methods;
    \item deterministic search over structured sample spaces.
\end{itemize}

This creates an important bridge between probabilistic combinatorics and
theoretical computer science.

%%%%%%%%%%%%%%%%%%%%%%%%%%%%%%%%%%%%%%%%%%%%%%%%%%%%%%%%%%%%
\subsection{Randomness Versus Pseudorandomness}
\label{subsec:randomness-pseudorandomness}
%%%%%%%%%%%%%%%%%%%%%%%%%%%%%%%%%%%%%%%%%%%%%%%%%%%%%%%%%%%%

A deterministic object may imitate many properties of a genuinely
random object.

Such structures are called pseudorandom.

Examples include graphs whose:

\begin{itemize}
    \item edge distribution resembles \(G(n,p)\);
    \item eigenvalues mimic random behavior;
    \item subgraph counts have expected densities;
    \item vertex subsets contain approximately expected numbers of
          edges.
\end{itemize}

Pseudorandomness is a major theme connecting algebra, probability,
graphs, and computation.

%%%%%%%%%%%%%%%%%%%%%%%%%%%%%%%%%%%%%%%%%%%%%%%%%%%%%%%%%%%%
\subsection{Chernoff Bounds Preview}
\label{subsec:chernoff-bounds-preview}
%%%%%%%%%%%%%%%%%%%%%%%%%%%%%%%%%%%%%%%%%%%%%%%%%%%%%%%%%%%%

Chebyshev's inequality provides polynomial tail bounds.

For sums of independent Bernoulli variables, much stronger exponential
bounds are available.

If

\[
X
=
\sum_{i=1}^{n}X_i
\]

with independent Bernoulli variables and

\[
\mu=\mathbb{E}[X],
\]

a typical Chernoff upper-tail bound has the form

\[
\boxed{
\mathbb{P}
\left(
X\geq(1+\delta)\mu
\right)
\leq
\exp
\left(
-\frac{\delta^2}{2+\delta}\mu
\right),
}
\]

for

\[
\delta>0.
\]

The Greek letter

\[
\delta
\]

is called ``delta.''

%%%%%%%%%%%%%%%%%%%%%%%%%%%%%%%%%%%%%%%%%%%%%%%%%%%%%%%%%%%%
\subsection{Why Concentration Matters}
\label{subsec:why-concentration-matters}
%%%%%%%%%%%%%%%%%%%%%%%%%%%%%%%%%%%%%%%%%%%%%%%%%%%%%%%%%%%%

Expectation describes the average.

Concentration tells us whether typical outcomes are close to that
average.

In many large random discrete structures,

\[
\boxed{
\text{random variable}
\approx
\text{expected value}
}
\]

with very high probability.

Concentration inequalities make this statement quantitative.

%%%%%%%%%%%%%%%%%%%%%%%%%%%%%%%%%%%%%%%%%%%%%%%%%%%%%%%%%%%%
\subsection{Jensen-Type and Convexity Ideas}
\label{subsec:convexity-probabilistic-combinatorics}
%%%%%%%%%%%%%%%%%%%%%%%%%%%%%%%%%%%%%%%%%%%%%%%%%%%%%%%%%%%%

Convexity inequalities can also produce probabilistic-combinatorial
bounds.

For a convex function \(\varphi\),

\[
\boxed{
\varphi(\mathbb{E}[X])
\leq
\mathbb{E}[\varphi(X)].
}
\]

This is Jensen's inequality.

Such inequalities can convert random degree distributions or random
counts into deterministic estimates.

%%%%%%%%%%%%%%%%%%%%%%%%%%%%%%%%%%%%%%%%%%%%%%%%%%%%%%%%%%%%
\subsection{Random Ordering Method}
\label{subsec:random-ordering-method}
%%%%%%%%%%%%%%%%%%%%%%%%%%%%%%%%%%%%%%%%%%%%%%%%%%%%%%%%%%%%

Another common strategy chooses a uniformly random ordering of a finite
set.

For example, one may:

\begin{itemize}
    \item place vertices in a random order;
    \item choose an object if it appears before all conflicting objects;
    \item calculate the probability that each object survives.
\end{itemize}

This technique frequently gives lower bounds for independent sets,
matchings, and packing problems.

%%%%%%%%%%%%%%%%%%%%%%%%%%%%%%%%%%%%%%%%%%%%%%%%%%%%%%%%%%%%
\subsection{Worked Example: Random Ordering and Independent Sets}
\label{subsec:worked-random-ordering-independent}
%%%%%%%%%%%%%%%%%%%%%%%%%%%%%%%%%%%%%%%%%%%%%%%%%%%%%%%%%%%%

\begin{workedexamplebox}{Choosing Vertices Before Their Neighbors}

Let \(G=(V,E)\).

Choose a uniformly random ordering of the vertices.

Include a vertex \(v\) if it appears before every vertex in its
neighborhood.

The selected vertices form an independent set.

Among

\[
\deg(v)+1
\]

vertices consisting of \(v\) and its neighbors, each is equally likely
to appear first.

Therefore,

\[
\mathbb{P}(v\text{ is selected})
=
\frac{1}{\deg(v)+1}.
\]

By linearity of expectation, the expected size of the independent set
is

\[
\sum_{v\in V}
\frac{1}{\deg(v)+1}.
\]

\begin{answerbox}
\[
\boxed{
\alpha(G)
\geq
\sum_{v\in V}
\frac{1}{\deg(v)+1}.
}
\]
\end{answerbox}

This is the Caro--Wei bound.

\end{workedexamplebox}

%%%%%%%%%%%%%%%%%%%%%%%%%%%%%%%%%%%%%%%%%%%%%%%%%%%%%%%%%%%%
\subsection{Random Greedy Methods}
\label{subsec:random-greedy-methods}
%%%%%%%%%%%%%%%%%%%%%%%%%%%%%%%%%%%%%%%%%%%%%%%%%%%%%%%%%%%%

A random greedy process repeatedly chooses an allowable object at random
and removes incompatible choices.

Examples include:

\begin{itemize}
    \item constructing matchings;
    \item building independent sets;
    \item generating sparse graphs;
    \item constructing packings;
    \item building combinatorial designs.
\end{itemize}

The process combines randomness with local optimization.

%%%%%%%%%%%%%%%%%%%%%%%%%%%%%%%%%%%%%%%%%%%%%%%%%%%%%%%%%%%%
\subsection{Entropy Method Preview}
\label{subsec:entropy-method-preview}
%%%%%%%%%%%%%%%%%%%%%%%%%%%%%%%%%%%%%%%%%%%%%%%%%%%%%%%%%%%%

Entropy provides another probabilistic tool for counting.

For a discrete random variable \(X\) with probabilities \(p_x\), its
Shannon entropy is

\[
\boxed{
H(X)
=
-\sum_x p_x\log p_x.
}
\]

Entropy measures uncertainty.

In combinatorics, entropy methods can prove upper bounds on the number
of possible configurations.

They are especially effective in:

\begin{itemize}
    \item set systems;
    \item graph enumeration;
    \item coding theory;
    \item information-theoretic inequalities.
\end{itemize}

%%%%%%%%%%%%%%%%%%%%%%%%%%%%%%%%%%%%%%%%%%%%%%%%%%%%%%%%%%%%
\subsection{Applications}
\label{subsec:probabilistic-combinatorics-applications}
%%%%%%%%%%%%%%%%%%%%%%%%%%%%%%%%%%%%%%%%%%%%%%%%%%%%%%%%%%%%

\begin{applicationbox}

\textbf{Extremal Combinatorics.}
Random constructions prove the existence of large structures avoiding
forbidden configurations.

\medskip

\textbf{Ramsey Theory.}
Random colorings provide strong lower bounds for Ramsey numbers.

\medskip

\textbf{Graph Theory.}
Random graphs model typical graph structure, threshold behavior,
subgraph counts, degree distributions, and connectivity.

\medskip

\textbf{Coding Theory.}
Random codes prove the existence of large collections of codewords with
prescribed minimum distance.

\medskip

\textbf{Computer Science.}
Randomized algorithms, derandomization, hashing, load balancing, and
random graph algorithms rely on probabilistic combinatorial methods.

\medskip

\textbf{Optimization.}
Randomized rounding converts fractional solutions into discrete
solutions with controlled expected cost and constraint violations.

\medskip

\textbf{Design Theory.}
Random and random-greedy constructions can produce packings and
approximate designs.

\medskip

\textbf{Network Science.}
Random graph models provide baseline models for large networks.

\medskip

\textbf{Statistical Physics.}
Random discrete configurations arise in percolation, spin systems, and
random constraint models.

\end{applicationbox}

%%%%%%%%%%%%%%%%%%%%%%%%%%%%%%%%%%%%%%%%%%%%%%%%%%%%%%%%%%%%
\subsection{Common Errors}
\label{subsec:probabilistic-combinatorics-common-errors}
%%%%%%%%%%%%%%%%%%%%%%%%%%%%%%%%%%%%%%%%%%%%%%%%%%%%%%%%%%%%

\subsubsection{Assuming Positive Expectation Implies Positive Value for Every Outcome}

If

\[
\mathbb{E}[X]>0,
\]

then some outcome has \(X>0\), but not every outcome must.

%%%%%%%%%%%%%%%%%%%%%%%%%%%%%%%%%%%%%%%%%%%%%%%%%%%%%%%%%%%%
\subsubsection{Assuming Independence Is Required for Linearity of Expectation}

Linearity gives

\[
\mathbb{E}[X+Y]
=
\mathbb{E}[X]
+
\mathbb{E}[Y]
\]

whether or not \(X\) and \(Y\) are independent.

%%%%%%%%%%%%%%%%%%%%%%%%%%%%%%%%%%%%%%%%%%%%%%%%%%%%%%%%%%%%
\subsubsection{Assuming the Union Bound Is Exact}

In general,

\[
\mathbb{P}(A\cup B)
\neq
\mathbb{P}(A)+\mathbb{P}(B).
\]

The union bound states only

\[
\mathbb{P}(A\cup B)
\leq
\mathbb{P}(A)+\mathbb{P}(B).
\]

%%%%%%%%%%%%%%%%%%%%%%%%%%%%%%%%%%%%%%%%%%%%%%%%%%%%%%%%%%%%
\subsubsection{Ignoring Dependence Between Complex Events}

Elementary choices may be independent while larger events are not.

Two subgraph events that share edges may be dependent.

%%%%%%%%%%%%%%%%%%%%%%%%%%%%%%%%%%%%%%%%%%%%%%%%%%%%%%%%%%%%
\subsubsection{Using Markov's Inequality for Variables That Can Be Negative}

Markov's inequality requires

\[
X\geq0.
\]

%%%%%%%%%%%%%%%%%%%%%%%%%%%%%%%%%%%%%%%%%%%%%%%%%%%%%%%%%%%%
\subsubsection{Confusing Expected Value with Typical Value}

A random variable can have a meaningful expectation but still vary
widely.

Variance and concentration determine whether the expectation is
representative.

%%%%%%%%%%%%%%%%%%%%%%%%%%%%%%%%%%%%%%%%%%%%%%%%%%%%%%%%%%%%
\subsubsection{Using the Local Lemma Without Checking Dependency}

The Local Lemma requires control over which bad events depend on one
another.

Small individual probabilities alone are not sufficient.

%%%%%%%%%%%%%%%%%%%%%%%%%%%%%%%%%%%%%%%%%%%%%%%%%%%%%%%%%%%%
\subsubsection{Assuming a Probabilistic Proof Gives an Efficient Construction}

A probabilistic existence proof may not immediately provide an
algorithm.

Derandomization is a separate question.

%%%%%%%%%%%%%%%%%%%%%%%%%%%%%%%%%%%%%%%%%%%%%%%%%%%%%%%%%%%%
\subsubsection{Confusing Threshold Scale with Exact Threshold}

An expectation calculation may suggest the correct scale

\[
p(n),
\]

but rigorous threshold behavior may require second moments,
concentration, or deeper methods.

%%%%%%%%%%%%%%%%%%%%%%%%%%%%%%%%%%%%%%%%%%%%%%%%%%%%%%%%%%%%
\subsection{Exercises}
\label{subsec:probabilistic-combinatorics-exercises}
%%%%%%%%%%%%%%%%%%%%%%%%%%%%%%%%%%%%%%%%%%%%%%%%%%%%%%%%%%%%

\begin{exercise}[1]
State the basic probabilistic method.
\end{exercise}

\begin{exercise}[1]
Let \(X\) be a nonnegative integer-valued random variable with

\[
\mathbb{E}[X]<1.
\]

Explain why there must be an outcome with \(X=0\).
\end{exercise}

\begin{exercise}[1]
Define an indicator random variable.
\end{exercise}

\begin{exercise}[1]
Show that

\[
\mathbb{E}[\mathbf{1}_A]
=
\mathbb{P}(A).
\]
\end{exercise}

\begin{exercise}[1]
State linearity of expectation.
\end{exercise}

\begin{exercise}[1]
Does linearity of expectation require independence?
\end{exercise}

\begin{exercise}[1]
In \(G(n,p)\), find the probability that a specified edge is present.
\end{exercise}

\begin{exercise}[1]
Find the expected number of edges in

\[
G(20,\tfrac12).
\]
\end{exercise}

\begin{exercise}[1]
Find the expected degree of a vertex in

\[
G(20,\tfrac12).
\]
\end{exercise}

\begin{exercise}[1]
Find the expected number of triangles in

\[
G(8,\tfrac12).
\]
\end{exercise}

\begin{exercise}[2]
Let \(R\subseteq[n]\) be formed by selecting each element independently
with probability \(p\). Prove

\[
\mathbb{E}[|R|]=np.
\]
\end{exercise}

\begin{exercise}[2]
Show that a uniformly random permutation of \(n\) elements has expected
number of fixed points equal to \(1\).
\end{exercise}

\begin{exercise}[2]
Find the expected number of inversions in a uniformly random
permutation of \(8\) elements.
\end{exercise}

\begin{exercise}[2]
For \(G\sim G(n,p)\), derive

\[
\mathbb{E}[|E(G)|]
=
p\binom n2.
\]
\end{exercise}

\begin{exercise}[2]
For \(G\sim G(n,p)\), derive

\[
\mathbb{E}[\#K_4]
=
\binom n4p^6.
\]
\end{exercise}

\begin{exercise}[2]
Find the expected number of independent \(k\)-sets in \(G(n,p)\).
\end{exercise}

\begin{exercise}[2]
State the union bound.
\end{exercise}

\begin{exercise}[2]
Suppose

\[
\sum_{i=1}^{m}\mathbb{P}(A_i)<1.
\]

Prove that there is an outcome avoiding every \(A_i\).
\end{exercise}

\begin{exercise}[2]
Use Markov's inequality when

\[
\mathbb{E}[X]=12
\]

to bound

\[
\mathbb{P}(X\geq60).
\]
\end{exercise}

\begin{exercise}[2]
Let

\[
X\sim\operatorname{Bin}(n,p).
\]

Find \(\mathbb{E}[X]\) and \(\operatorname{Var}(X)\).
\end{exercise}

\begin{exercise}[2]
State Chebyshev's inequality.
\end{exercise}

\begin{exercise}[2]
Suppose \(X\) has mean \(100\) and variance \(25\). Use Chebyshev's
inequality to bound

\[
\mathbb{P}(|X-100|\geq10).
\]
\end{exercise}

\begin{exercise}[3]
Use the probabilistic method to prove that some cut in every graph with
\(m\) edges contains at least

\[
\frac m2
\]

edges.
\end{exercise}

\begin{exercise}[3]
Hint for the previous exercise: independently place each vertex into
one of two parts with equal probability.
\end{exercise}

\begin{exercise}[3]
Let \(G\) have \(n\) vertices and \(m\) edges. Derive the alteration
bound

\[
\alpha(G)
\geq
pn-p^2m.
\]
\end{exercise}

\begin{exercise}[3]
Optimize the previous bound over \(p\).
\end{exercise}

\begin{exercise}[3]
Prove the Caro--Wei bound

\[
\alpha(G)
\geq
\sum_{v\in V}
\frac{1}{\deg(v)+1}
\]

using a random ordering of the vertices.
\end{exercise}

\begin{exercise}[3]
Use random two-colorings of the edges of \(K_n\) to prove that if

\[
\binom nk
2^{1-\binom{k}{2}}
<1,
\]

then there exists a coloring with no monochromatic \(K_k\).
\end{exercise}

\begin{exercise}[3]
Use the previous result to derive a lower bound for the diagonal Ramsey
number \(R(k,k)\).
\end{exercise}

\begin{exercise}[3]
For \(G\sim G(n,p)\), compute the expected number of copies of a fixed
tree on \(k\) vertices.
\end{exercise}

\begin{exercise}[3]
For \(G\sim G(n,p)\), compute the expected number of cycles of length
\(r\).
\end{exercise}

\begin{exercise}[3]
Use a first moment argument to identify the natural threshold scale for
the appearance of a triangle in \(G(n,p)\).
\end{exercise}

\begin{exercise}[3]
Explain why a first moment calculation alone does not always prove that
a structure exists with high probability.
\end{exercise}

\begin{exercise}[3]
Give an example of two dependent triangle events in \(G(n,p)\).
\end{exercise}

\begin{exercise}[3]
State the symmetric Lov\'asz Local Lemma.
\end{exercise}

\begin{exercise}[3]
Explain why the Local Lemma can succeed even when the union bound is
larger than \(1\).
\end{exercise}

\begin{exercise}[3]
Describe how conditional expectation can convert an existence proof
into a deterministic algorithm.
\end{exercise}

\begin{exercise}[4]
Study the second moment method and derive

\[
\mathbb{P}(X>0)
\geq
\frac{\mathbb{E}[X]^2}{\mathbb{E}[X^2]}
\]

for a nonnegative random variable \(X\).
\end{exercise}

\begin{exercise}[4]
Investigate Chernoff bounds for sums of independent Bernoulli random
variables.
\end{exercise}

\begin{exercise}[4]
Study threshold functions for connectivity in \(G(n,p)\).
\end{exercise}

\begin{exercise}[4]
Investigate the threshold for the appearance of \(K_r\) in \(G(n,p)\).
\end{exercise}

\begin{exercise}[4]
Study the asymptotic size of the largest clique in \(G(n,\tfrac12)\).
\end{exercise}

\begin{exercise}[4]
Study the asymptotic size of the largest independent set in
\(G(n,\tfrac12)\).
\end{exercise}

\begin{exercise}[4]
Investigate the Lov\'asz Local Lemma through a graph-coloring
application.
\end{exercise}

\begin{exercise}[4]
Study the Moser--Tardos algorithmic version of the Local Lemma.
\end{exercise}

\begin{exercise}[4]
Investigate randomized rounding in combinatorial optimization.
\end{exercise}

\begin{exercise}[4]
Study the entropy method in combinatorics and prove one elementary
counting inequality using entropy.
\end{exercise}

%%%%%%%%%%%%%%%%%%%%%%%%%%%%%%%%%%%%%%%%%%%%%%%%%%%%%%%%%%%%
\subsection{Conceptual Summary}
\label{subsec:probabilistic-combinatorics-summary}
%%%%%%%%%%%%%%%%%%%%%%%%%%%%%%%%%%%%%%%%%%%%%%%%%%%%%%%%%%%%

Probabilistic combinatorics uses randomness to study deterministic
discrete structures.

The basic probabilistic method is

\[
\boxed{
\mathbb{P}(\text{desired property})>0
\Longrightarrow
\text{a desired object exists}.
}
\]

Indicator variables satisfy

\[
\boxed{
\mathbb{E}[\mathbf{1}_A]
=
\mathbb{P}(A).
}
\]

Linearity of expectation gives

\[
\boxed{
\mathbb{E}
\left[
\sum_iX_i
\right]
=
\sum_i\mathbb{E}[X_i].
}
\]

The first moment method often uses

\[
\boxed{
\mathbb{E}[X]<1
}
\]

to prove the existence of an outcome with \(X=0\).

The union bound states

\[
\boxed{
\mathbb{P}
\left(
\bigcup_iA_i
\right)
\leq
\sum_i\mathbb{P}(A_i).
}
\]

The Erd\H{o}s--R\'enyi model

\[
\boxed{
G(n,p)
}
\]

includes each possible edge independently with probability \(p\).

Its expected number of edges is

\[
\boxed{
p\binom n2.
}
\]

Its expected number of \(K_k\)'s is

\[
\boxed{
\binom nk
p^{\binom{k}{2}}.
}
\]

Variance and concentration inequalities measure fluctuations around
expectation.

Markov's inequality gives

\[
\boxed{
\mathbb{P}(X\geq a)
\leq
\frac{\mathbb{E}[X]}{a}
}
\]

for \(X\geq0\).

Chebyshev's inequality gives

\[
\boxed{
\mathbb{P}
\left(
|X-\mathbb{E}[X]|\geq t
\right)
\leq
\frac{\operatorname{Var}(X)}{t^2}.
}
\]

The Lov\'asz Local Lemma shows that rare, locally dependent bad events
can sometimes all be avoided simultaneously.

The alteration method uses

\[
\boxed{
\text{random construction}
+
\text{deterministic repair}.
}
\]

Random graph theory studies thresholds, typical structure, and
concentration.

Randomized methods can also be derandomized through conditional
expectation and related techniques.

The broad principle is

\[
\boxed{
\text{random model}
\longrightarrow
\text{probability or expectation}
\longrightarrow
\text{deterministic combinatorial conclusion}.
}
\]

%%%%%%%%%%%%%%%%%%%%%%%%%%%%%%%%%%%%%%%%%%%%%%%%%%%%%%%%%%%%
\subsection{Section Connections}
\label{subsec:probabilistic-combinatorics-connections}
%%%%%%%%%%%%%%%%%%%%%%%%%%%%%%%%%%%%%%%%%%%%%%%%%%%%%%%%%%%%

\chapterconnection{
Probabilistic Combinatorics builds on the existence questions of
Extremal Combinatorics and introduces randomness as a systematic proof
tool. The next section, Graph Theory, develops deterministic graph
structure in depth, while random graphs provide an important parallel
model for understanding typical graphs. Ramsey Theory uses random
colorings to obtain lower bounds for Ramsey numbers. Design Theory and
Matroid Theory use random and randomized constructions in existence and
optimization problems. Chapter~\ref{chap:probability} develops the
probability theory underlying expectation, variance, concentration, and
random variables in substantially greater depth. Optimization and
Operations Research connects through randomized rounding and
probabilistic algorithms.
}

%%%%%%%%%%%%%%%%%%%%%%%%%%%%%%%%%%%%%%%%%%%%%%%%%%%%%%%%%%%%
% SECTION SOURCE TRACKING
%%%%%%%%%%%%%%%%%%%%%%%%%%%%%%%%%%%%%%%%%%%%%%%%%%%%%%%%%%%%

%------------------------------------------------------------
% 6.8 Probabilistic Combinatorics
%------------------------------------------------------------
%
% Source basis:
%   - Standard probabilistic combinatorics.
%   - Standard probabilistic method.
%   - Standard random graph theory.
%   - Standard introductory concentration methods.
%
% Used for:
%   - probabilistic method;
%   - existence from positive probability;
%   - indicator random variables;
%   - linearity of expectation;
%   - random subsets;
%   - first moment method;
%   - expected fixed points;
%   - Erdős--Rényi random graphs;
%   - expected edge counts;
%   - expected degrees;
%   - expected triangle counts;
%   - fixed-subgraph counts;
%   - random colorings;
%   - union bound;
%   - probabilistic Ramsey bounds;
%   - Markov's inequality;
%   - variance;
%   - Chebyshev's inequality;
%   - second moment method;
%   - independence and dependence;
%   - Lovász Local Lemma;
%   - randomized coloring;
%   - alteration method;
%   - independent-set bounds;
%   - random sampling and deletion;
%   - clique and independent-set expectations;
%   - threshold phenomena;
%   - with-high-probability terminology;
%   - random hypergraphs;
%   - random set systems;
%   - random permutations;
%   - expected inversion counts;
%   - expected cycle counts;
%   - randomized algorithms;
%   - conditional expectation;
%   - derandomization;
%   - pseudorandomness;
%   - Chernoff-bound preview;
%   - random ordering methods;
%   - Caro--Wei bound;
%   - random greedy processes;
%   - entropy-method preview;
%   - applications and exercises.
%
% Notes:
%   - Probability theory is developed systematically in Chapter 7.
%   - Random graph structure is introduced here only to support
%     combinatorial methods; deterministic graph theory appears next.
%   - Ramsey applications are developed further in the Ramsey Theory
%     section.
%   - Advanced concentration, martingale, entropy, and probabilistic
%     inequalities belong primarily to the Probability chapter.
%
%%%%%%%%%%%%%%%%%%%%%%%%%%%%%%%%%%%%%%%%%%%%%%%%%%%%%%%%%%%%